%%%%%%%%%%%%%%%%%%%%%%%%%%%%%%%%%%%%%%%%%%%%%%%%%%%%%%%%%%%%%%%%%%%%%%%%%%%%%%%
% Appendix I: Joint \ac{PDF}, Likelihood, and Factor Graph (Extended-Object Model)
%%%%%%%%%%%%%%%%%%%%%%%%%%%%%%%%%%%%%%%%%%%%%%%%%%%%%%%%%%%%%%%%%%%%%%%%%%%%%%%

\section{\bl{Derivation of the }\fg{Approximate }\bl{Joint PDF}}
\label{sec:appendix_I}

In this appendix, we derive the single-step joint \ac{PDF} and its factorization for the single-\ac{BS} filter of \secref{sec:single_bs_method}, \fg{conditioned on the externally produced partition, which we treat as given side information. Approximations are marked by \(\approx\): the clustering-induced group count is modeled as Poisson, so the resulting factorization is an approximate conditional model rather than a generative law over the raw measurements.} The derivation follows the point-target message-passing filters \cite{meyer2018messagepassing,meyer2018messagepassingSupplement,meyer2021scalable}, with two differences specific to our extended-object model: measurements are processed in groups rather than one per target, and the number of clutter groups produced by the clustering step is approximated by a Poisson distribution. The key idea is to express the unknown numbers of detected legacy targets, newborn targets, and clutter groups as deterministic functions of the association and existence variables, after which the joint \ac{PDF} factorizes into per-target terms. We omit the time index \(k\) throughout.

\newcommand{\Pm}[2]{{}_{#1}P_{#2}}  % permutation symbol
\newcommand{\mucgroup}{\mu_{g}}

\subsection*{A. Notation and Setup}

We reuse the notation of \secref{sec:system_model} and \secref{sec:single_bs_method}. The \(m\) measurements are grouped into \(n^g\) groups \(\V{Z}^{G}=\{\V{Z}^1,\dots,\V{Z}^{n^g}\}\), where group \(\V{Z}^g\) contains \(|\V{Z}^g|\) measurements and both \(m\) and \(n^g\) are determined by \(\V{Z}^{G}\); we use ``group'' rather than ``cluster'' to reserve the superscript \(c\) for clutter. The \(n^t\) legacy targets have augmented states \(\underline{\V{y}}^i=[\underline{\V{x}}^{i\T},\underline{\M{E}}^{i\T},\underline{r}^i]^\T\) stacked in \(\underline{\V{Y}}\). We introduce one newborn component per group, \(\overline{\V{y}}^j=[\overline{\V{x}}^{j\T},\overline{\M{E}}^{j\T},\overline{r}^j]^\T\) stacked in \(\overline{\V{Y}}\), write \(\overline{\Set{X}}=((\overline{\V{x}}^1,\overline{\M{E}}^1),\dots,(\overline{\V{x}}^{n^g},\overline{\M{E}}^{n^g}))\) for the newborn kinematic--extent pairs, and collect the legacy and newborn existence indicators in \(\underline{\V{r}}\) and \(\overline{\V{r}}\).

Assuming, as in \cite[Eq.~(36)]{meyer2018messagepassing} and \cite[Eqs.~(1)--(4)]{meyer2021scalable}, that the time-\((k{-}1)\) posterior and the transition density factorize into single-target terms, the legacy prior is
\begin{equation}
f(\underline{\V{Y}}) = \prod_{i=1}^{n^t} f(\underline{\V{y}}^i),
\label{eq:pt_legacy_prior}
\end{equation}
with \(f(\underline{\V{y}}^i)\) the predicted marginal of legacy target \(i\).

Following \cite[Eqs.~(10),(20)]{meyer2018messagepassing}, the target- and measurement-oriented association vectors are \(\V{a}=(a^1,\dots,a^{n^t})\) and \(\V{b}=(b^1,\dots,b^{n^g})\): \(a^i=g\) if legacy target \(i\) generates group \(g\) and \(a^i=0\) otherwise, while \(b^j=\ell\) if group \(j\) is generated by legacy target \(\ell\) and \(b^j=0\) otherwise. The consistency factor \(\boldsymbol{\psi}(\V{a},\V{b})=\prod_{i,j}\Psi_{i,j}(a^i,b^j)\), with \(\Psi_{i,j}=0\) iff exactly one of \(a^i=j\) and \(b^j=i\) holds and \(\Psi_{i,j}=1\) otherwise, restricts the model to exclusive association hypotheses; its marginals reproduce the single-sided exclusivity factors, \(\sum_{\V{b}}\boldsymbol{\psi}(\V{a},\V{b})=\boldsymbol{\psi}(\V{a})\) and \(\sum_{\V{a}}\boldsymbol{\psi}(\V{a},\V{b})=\boldsymbol{\psi}(\V{b})\).

Each group originates from a single source. We collect the indices of detected legacy targets, existing newborns, and clutter groups in \(\mathbb{D}(\V{a})=\{i:a^i\neq0\}\), \(\mathbb{N}(\overline{\V{r}})=\{j:\overline{r}^j=1\}\), and \(\mathbb{C}(\V{a},\overline{\V{r}})=\{1,\dots,n^g\}\setminus(\{a^i:i\in\mathbb{D}(\V{a})\}\cup\mathbb{N}(\overline{\V{r}}))\), with cardinalities \(n^d=|\mathbb{D}(\V{a})|\), \(n^n=|\mathbb{N}(\overline{\V{r}})|\), and \(n^c=|\mathbb{C}(\V{a},\overline{\V{r}})|=n^g-n^d-n^n\). These three counts are deterministic functions of \((\V{a},\overline{\V{r}})\).

\subsection*{B. Joint \ac{PDF} and its Factorization}

By the chain rule, and inserting the exclusivity factor \(\boldsymbol{\psi}(\V{a},\V{b})\), the joint \ac{PDF} decomposes into the legacy prior, an association model, and a likelihood conditioned on the association,\footnote{\fg{Here and throughout this appendix, \(f(\cdot,\mathbf Z^G)\) denotes the density conditional on the fixed partition \(\mathcal P=\mathcal C(\mathbf Z)\) supplied by the external clustering algorithm; dependence on \(\mathcal P\) is suppressed for readability. Equation (A.2) is therefore an exact chain-rule decomposition of this conditional density. The subsequent Poisson group-count and conditional group-independence models approximate its factors; we do not claim to model the probability that the clustering algorithm produces \(\mathcal P\). Moreover, the conditional model assumes that each target generates at most one group and each group contains measurements from at most one target; clustering splits and cross-target merges are not represented.}}
\begin{equation}
\begin{aligned}
f(\underline{\V{Y}},\overline{\V{Y}},\V{a},\V{b},\V{Z}^{G})
={}& f(\underline{\V{Y}})\,
\boldsymbol{\psi}(\V{a},\V{b})\,
f(\V{a},\V{b},\overline{\V{Y}}\mid\underline{\V{Y}})\\
&\times f(\V{Z}^{G}\mid\underline{\V{Y}},\overline{\V{Y}},\V{a},\V{b}).
\end{aligned}
\label{eq:joint_pdf_structured}
\end{equation}
The legacy prior is \eqref{eq:pt_legacy_prior}; we factorize the association model and the likelihood in turn, omitting \(\boldsymbol{\psi}(\V{a},\V{b})\) until the final result.

\emph{Association model.} We split the association model into an existence-and-association part and the newborn densities, \(f(\V{a},\V{b},\overline{\V{Y}}\mid\underline{\V{Y}})=f(\overline{\V{r}},\V{a},\V{b}\mid\underline{\V{Y}})\,f(\overline{\Set{X}}\mid\overline{\V{r}})\). The numbers of clutter groups and detected newborns are modeled as Poisson \cite[Eqs.~(3),(18)]{meyer2018messagepassingSupplement},
\begin{equation}
f(n^c)\approx\frac{\mucgroup^{\,n^c}\mathrm e^{-\mucgroup}}{n^c!},
\qquad
f(n^n)=\frac{\mu_n^{\,n^n}\mathrm e^{-\mu_n}}{n^n!},
\label{eq:pt_counts_distributions}
\end{equation}
where \(\mu_n\) is the mean number of detected newborns. The clustering step turns the clutter \ac{PPP} of mean \(\mu_c\) into groups; since one group may contain several clutter detections, \(\mucgroup\neq\mu_c\) in general, and we approximate the induced group count as Poisson with mean \(\mucgroup\). Conditioned on the counts, the newborn existence pattern is uniform, \(f(\overline{\V{r}}\mid\cdot)=\binom{n^c+n^n}{n^n}^{-1}\), and the associations are uniform over the \(\Pm{n^g}{n^d}\) feasible assignments, \(f(\V{a},\V{b}\mid\cdot)=1/\Pm{n^g}{n^d}=(n^c+n^n)!/n^g!\) \cite[Eqs.~(22)--(24)]{meyer2018messagepassingSupplement}. A legacy target contributes \(\underline{r}^p\) when detected and \(1-\underline{r}^p(1-\mathrm e^{-\mu_m})\) when not, where \(\mu_m=\mu_m(\V{x},\M{E})\) is the expected number of target measurements \eqref{eq:system_measurement_count_mean}; the newborn densities are \(\prod_{j\in\mathbb{N}(\overline{\V{r}})}f_n(\overline{\V{x}}^j,\overline{\M{E}}^j)\prod_{j\notin\mathbb{N}(\overline{\V{r}})}f_d(\overline{\V{x}}^j,\overline{\M{E}}^j)\), with \(f_n\) the newborn density and \(f_d\) a dummy density that integrates to one. Multiplying these factors, the \(n^c!\), \(n^n!\), and \((n^c+n^n)!\) terms cancel, and the association model %is \(\boldsymbol{\psi}(\V{a},\V{b})\) times
\begin{equation}
\begin{aligned}
&f(\V{a},\V{b},\overline{\V{Y}}\mid\underline{\V{Y}})
\approx \frac{\mucgroup^{\,n^c}\mathrm e^{-\mucgroup}\,\mu_n^{\,n^n}\mathrm e^{-\mu_n}}{n^g!}
\\
&\quad\times
\prod_{p\in\mathbb{D}(\V{a})}\underline{r}^{p}
\prod_{p\notin\mathbb{D}(\V{a})}\bigl[1-\underline{r}^{p}(1-\mathrm e^{-\mu_m})\bigr]
\\
&\quad\times
\prod_{j\in\mathbb{N}(\overline{\V{r}})}f_n(\overline{\V{x}}^{j},\overline{\M{E}}^{j})
\prod_{j\notin\mathbb{N}(\overline{\V{r}})}f_d(\overline{\V{x}}^{j},\overline{\M{E}}^{j}).
\end{aligned}
\label{eq:I13}
\end{equation}

\emph{Likelihood.} Conditioned on the association, the groups are generated independently,
\begin{equation}
\begin{aligned}
&f(\V{Z}^{G}\mid\underline{\V{Y}},\overline{\V{Y}},\V{a},\V{b})\\
&\quad=\prod_{p\in\mathbb{D}(\V{a})} f(\V{Z}^{a^{p}}\mid\underline{\V{x}}^{p},\underline{\M{E}}^{p})\\
&\qquad\times\prod_{j\in\mathbb{N}(\overline{\V{r}})} f(\V{Z}^{j}\mid\overline{\V{x}}^{j},\overline{\M{E}}^{j})
\prod_{c\in\mathbb{C}(\V{a},\overline{\V{r}})} f_c(\V{Z}^{c}).
\end{aligned}
\label{eq:pt_group_like_count_split}
\end{equation}
A detected legacy target follows the standard extended-target \ac{PPP} model \cite[Eq.~(24)]{Yuxuanxia_standardmodel}, \(f(\V{Z}\mid\V{x},\M{E})=\mathrm e^{-\mu_m}\prod_{\V{z}\in\V{Z}}\mu_m f(\V{z}\mid\V{x},\M{E})\). A newborn is detected by assumption, so its likelihood is the zero-truncated version \(f(\V{Z}\mid\V{x},\M{E})=[\mathrm e^{-\mu_m}/(1-\mathrm e^{-\mu_m})]\prod_{\V{z}\in\V{Z}}\mu_m f(\V{z}\mid\V{x},\M{E})\). A clutter group is i.i.d., \(f_c(\V{Z}^{c})=\prod_{\V{z}\in\V{Z}^{c}}f_c(\V{z})\), and since \(\mathbb{D}\), \(\mathbb{N}\), and \(\mathbb{C}\) partition the groups, \(\prod_{c=1}^{n^g}f_c(\V{Z}^{c})=\prod_{i=1}^{m}f_c(\V{z}^{i})\).

\emph{Factorized joint \ac{PDF}.} Combining the prior \eqref{eq:pt_legacy_prior}, the association model \eqref{eq:I13}, and the likelihood \eqref{eq:pt_group_like_count_split}, dividing each detected group by \(f_c\), and redistributing \(\mucgroup^{\,n^c}=\mucgroup^{\,n^g}\prod_{p\in\mathbb{D}(\V{a})}\mucgroup^{-1}\prod_{j\in\mathbb{N}(\overline{\V{r}})}\mucgroup^{-1}\), all terms independent of the unknowns collect into the constant
\begin{equation}
C(\mucgroup,\mu_n,\V{Z}^{G})
\triangleq
\frac{\mucgroup^{\,n^g}\,\mathrm e^{-\mucgroup}\,\mathrm e^{-\mu_n}}{n^g!}
\prod_{i=1}^{m} f_c(\V{z}^{i}),
\label{eq:pt_constant_C}
\end{equation}
where \(n^g\) and \(m\) are determined by \(\V{Z}^{G}\). To lighten the per-target factors, we suppress the component state, writing \(\mu_m\) for \(\mu_m(\V{x},\M{E})\), \(f(\V{z})\) for \(f(\V{z}\mid\V{x},\M{E})\), and \(f_n,f_d\) for \(f_n(\overline{\V{x}}^{j},\overline{\M{E}}^{j}),f_d(\overline{\V{x}}^{j},\overline{\M{E}}^{j})\). The legacy likelihood factor is
\begin{equation}
\underline{l}(\underline{\V{y}}^{p},a^{p};\V{Z}^{G})
=
\begin{cases}
\begin{aligned}[t]
&\dfrac{\mathrm e^{-\mu_m}}{\mucgroup}\\
&\times\prod_{\V{z}\in\V{Z}^{a^{p}}}\dfrac{\mu_m f(\V{z})}{f_c(\V{z})},
\end{aligned}
& a^{p}\neq 0,\ \underline{r}^{p}=1,\\[1.2ex]
\mathrm e^{-\mu_m}, & a^{p}=0,\ \underline{r}^{p}=1,\\[0.6ex]
1, & a^{p}=0,\ \underline{r}^{p}=0,\\[0.6ex]
0, & a^{p}\neq 0,\ \underline{r}^{p}=0,
\end{cases}
\label{eq:legacy_l_function}
\end{equation}
and the newborn likelihood factor is
\begin{equation}
\overline{l}(\overline{\V{y}}^{j},b^{j};\V{Z}^{j})
=
\begin{cases}
\begin{aligned}[t]
&\dfrac{\mu_n f_n\,\mathrm e^{-\mu_m}}{\mucgroup(1-\mathrm e^{-\mu_m})}\\
&\times\prod_{\V{z}\in\V{Z}^{j}}\dfrac{\mu_m f(\V{z})}{f_c(\V{z})},
\end{aligned}
& b^{j}=0,\ \overline{r}^{j}=1,\\[1.2ex]
0, & b^{j}\neq 0,\ \overline{r}^{j}=1,\\[0.6ex]
f_d, & \overline{r}^{j}=0.
\end{cases}
\label{eq:new_l_function}
\end{equation}
With these definitions, the joint \ac{PDF} factorizes into per-target terms as
\begin{equation}
\begin{aligned}
f(\underline{\V{Y}},\overline{\V{Y}},\V{a},\V{b},\V{Z}^{G})
\approx{}& C(\mucgroup,\mu_n,\V{Z}^{G})\,\boldsymbol{\psi}(\V{a},\V{b})\\
&\times \prod_{p=1}^{n^t}\bigl[f(\underline{\V{y}}^{p})\,\underline{l}(\underline{\V{y}}^{p},a^{p};\V{Z}^{G})\bigr]\\
&\times \prod_{j=1}^{n^{g}}\overline{l}(\overline{\V{y}}^{j},b^{j};\V{Z}^{j}),
\end{aligned}
\label{eq:final_joint_pdf_local_factors}
\end{equation}
which is the per-target product used in \secref{sec:single_bs_method}.

